\documentclass[11pt]{article}
\usepackage[utf8]{inputenc}	% Para caracteres en español
\usepackage{amsmath,amsthm,amsfonts,amssymb,amscd}
\usepackage{multirow,booktabs}
\usepackage[table]{xcolor}
\usepackage{fullpage}
\usepackage{lastpage}
\usepackage{enumitem}
\usepackage{fancyhdr}
\usepackage{mathrsfs}
\usepackage{wrapfig}
\usepackage{setspace}
\usepackage{calc}
\usepackage{multicol}
\usepackage{cancel}
\usepackage[retainorgcmds]{IEEEtrantools}
\usepackage[margin=3cm]{geometry}
\usepackage{amsmath}
\newlength{\tabcont}
\setlength{\parindent}{0.0in}
\setlength{\parskip}{0.05in}
\usepackage{empheq}
\usepackage{framed}
\usepackage[most]{tcolorbox}
\usepackage{xcolor}
\colorlet{shadecolor}{orange!15}
\parindent 0in
\parskip 12pt
\geometry{margin=1in, headsep=0.25in}
\theoremstyle{definition}
\newtheorem{defn}{Definition}
\newtheorem{reg}{Rule}
\newtheorem{exer}{Exercise}
\newtheorem{note}{Note}
\begin{document}
\setcounter{section}{0}
\title{Econ 457 Lecture Notes}

\thispagestyle{empty}

% Some tools:
% Equation: \begin{equation}
% Notes: \begin{note}
% Shaded area: \begin{shaded}
% Definition: \begin{defn}


\begin{center}
{\LARGE \bf Econ 357}\\
{\bf Lecture Notes}\\
\end{center}

\section{Lecture 1}
\subsection{Balance Sheets}

\begin{itemize}
\item \textbf{Balance Sheets} 
\end{itemize}

\centering
\includegraphics[width=1\textwidth]{Econ357/figures/Notes_hh_bs.png}


\centering
\includegraphics[width=1\textwidth]{Econ357/figures/Notes_corp_bs.png}



\end{document}